\documentclass[11pt]{article}

\usepackage[margin=0.75in]{geometry}
\usepackage[T1]{fontenc}
\usepackage[utf8]{inputenc}
\usepackage{lmodern}
\usepackage{amsmath, amssymb}
\usepackage{xcolor}
\usepackage{enumitem}
\usepackage{hyperref}
\usepackage{tikz}
\usepackage{pgfplots}
\usepgfplotslibrary{fillbetween}
\pgfplotsset{compat=1.18}

\hypersetup{
	colorlinks=true,
	linkcolor=red!60!black,
	urlcolor=blue!60!black
}

\setlength{\parindent}{0pt}
\setlength{\parskip}{6pt}

\newcommand{\StandardNormalLeftTailGraph}[3]{%
	\begin{center}
		\begin{tikzpicture}
			\pgfmathsetmacro{\zbound}{#1}
			\pgfmathsetmacro{\ybound}{0.3989422804*exp(-(\zbound*\zbound)/2)}
			\begin{axis}[
				width=0.78\textwidth,
				height=5.4cm,
				xmin=-3.5,
				xmax=3.5,
				ymin=0,
				ymax=0.45,
				axis lines=left,
				xlabel={$z$},
				ylabel={$\phi(z)$},
				grid=both,
				xtick={-3,-2,-1,0,1,2,3},
				ytick={0,0.1,0.2,0.3,0.4},
				tick label style={font=\small},
				label style={font=\small},
				every axis plot/.append style={thick},
				clip=false
				]
				\addplot[name path=xaxis, draw=none] coordinates {(-3.5,0) (3.5,0)};
				\addplot[name path=curve, black, domain=-3.5:3.5, samples=220] {0.3989422804*exp(-x^2/2)};
				\addplot[name path=shadepath, draw=none, domain=-3.5:\zbound, samples=220] {0.3989422804*exp(-x^2/2)};
				\addplot[blue!35] fill between[of=shadepath and xaxis];
				\addplot[black, dashed] coordinates {(\zbound,0) (\zbound,\ybound)};
				\node[below, font=\small] at (axis cs:\zbound,0) {$#2$};
				\draw[->, thick] (axis cs:-2.05,0.31) -- (axis cs:-2.80,0.08);
				\node[font=\small, fill=white, inner sep=1.5pt] at (axis cs:-1.70,0.33) {$#3$};
			\end{axis}
		\end{tikzpicture}
	\end{center}
}

\newcommand{\StandardNormalRightTailGraph}[3]{%
	\begin{center}
		\begin{tikzpicture}
			\pgfmathsetmacro{\zbound}{#1}
			\pgfmathsetmacro{\ybound}{0.3989422804*exp(-(\zbound*\zbound)/2)}
			\begin{axis}[
				width=0.78\textwidth,
				height=5.4cm,
				xmin=-3.5,
				xmax=3.5,
				ymin=0,
				ymax=0.45,
				axis lines=left,
				xlabel={$z$},
				ylabel={$\phi(z)$},
				grid=both,
				xtick={-3,-2,-1,0,1,2,3},
				ytick={0,0.1,0.2,0.3,0.4},
				tick label style={font=\small},
				label style={font=\small},
				every axis plot/.append style={thick},
				clip=false
				]
				\addplot[name path=xaxis, draw=none] coordinates {(-3.5,0) (3.5,0)};
				\addplot[name path=curve, black, domain=-3.5:3.5, samples=220] {0.3989422804*exp(-x^2/2)};
				\addplot[name path=shadepath, draw=none, domain=\zbound:3.5, samples=220] {0.3989422804*exp(-x^2/2)};
				\addplot[orange!45] fill between[of=shadepath and xaxis];
				\addplot[black, dashed] coordinates {(\zbound,0) (\zbound,\ybound)};
				\node[below, font=\small] at (axis cs:\zbound,0) {$#2$};
				\draw[->, thick] (axis cs:2.00,0.31) -- (axis cs:2.75,0.08);
				\node[font=\small, fill=white, inner sep=1.5pt] at (axis cs:1.65,0.33) {$#3$};
			\end{axis}
		\end{tikzpicture}
	\end{center}
}

\newcommand{\StandardNormalCentralBandGraph}[5]{%
	\begin{center}
		\begin{tikzpicture}
			\pgfmathsetmacro{\zleft}{#1}
			\pgfmathsetmacro{\zright}{#2}
			\pgfmathsetmacro{\yleft}{0.3989422804*exp(-(\zleft*\zleft)/2)}
			\pgfmathsetmacro{\yright}{0.3989422804*exp(-(\zright*\zright)/2)}
			\begin{axis}[
				width=0.78\textwidth,
				height=5.4cm,
				xmin=-3.5,
				xmax=3.5,
				ymin=0,
				ymax=0.45,
				axis lines=left,
				xlabel={$z$},
				ylabel={$\phi(z)$},
				grid=both,
				xtick={-3,-2,-1,0,1,2,3},
				ytick={0,0.1,0.2,0.3,0.4},
				tick label style={font=\small},
				label style={font=\small},
				every axis plot/.append style={thick},
				clip=false
				]
				\addplot[name path=xaxis, draw=none] coordinates {(-3.5,0) (3.5,0)};
				\addplot[name path=curve, black, domain=-3.5:3.5, samples=220] {0.3989422804*exp(-x^2/2)};
				\addplot[name path=shadepath, draw=none, domain=\zleft:\zright, samples=220] {0.3989422804*exp(-x^2/2)};
				\addplot[blue!35] fill between[of=shadepath and xaxis];
				\addplot[black, dashed] coordinates {(\zleft,0) (\zleft,\yleft)};
				\addplot[black, dashed] coordinates {(\zright,0) (\zright,\yright)};
				\node[below, font=\small] at (axis cs:\zleft,0) {$#3$};
				\node[below, font=\small] at (axis cs:\zright,0) {$#4$};
				\draw[->, thick] (axis cs:0,0.34) -- (axis cs:0,0.17);
				\node[font=\small, fill=white, inner sep=1.5pt] at (axis cs:0,0.37) {$#5$};
			\end{axis}
		\end{tikzpicture}
	\end{center}
}

\begin{document}
	
	\begin{center}
		{\LARGE Term 3 Practice Activity: C7 Inverse Normal Distribution}
	\end{center}
	
	\tableofcontents
	
	\section*{Evaluated criteria}
	\begin{itemize}[leftmargin=*]
		\item C7.
	\end{itemize}
	
	\section{Top 5\% Exam Scores}
	
	\textbf{Problem.} A large exam has scores modeled by
	\[
	X \sim N(72, 9^2).
	\]
	
	\textbf{Questions/Tasks.} Identify the target probability before using inverse normal, then find the score $x$ that marks the top 5\% of students.
	
	\subsection*{C7}
	
	Top 5\% means
	\[
	P(X > x) = 0.05 \Rightarrow P(X \leq x) = 0.95.
	\]
	
	\[
	P(Z>1.645)=0.05.
	\]
	
	\StandardNormalRightTailGraph{1.645}{z=1.645}{P(Z>1.645)=0.05}
	
	Standardize:
	\[
	P\left(Z \leq \frac{x - 72}{9}\right) = 0.95.
	\]
	
	From the inverse normal table, $z_{0.95} \approx 1.645$. So
	\[
	\frac{x - 72}{9} = 1.645 \Rightarrow x = 72 + 1.645(9) = 86.805.
	\]
	
	\[
	x \approx 86.8
	\]
	
	\textbf{Interpretation.} A score of about 86.8 is the cutoff so only about 5\% of students score higher.
	
	\section{Bottom 10\% Commute Times}
	
	\textbf{Problem.} Morning commute times (minutes) are modeled by
	\[
	X \sim N(38, 7^2).
	\]
	
	\textbf{Questions/Tasks.} Identify the target probability before using inverse normal, then find the time $x$ below which the bottom 10\% of commutes fall.
	
	\subsection*{C7}
	
	Bottom 10\% means
	\[
	P(X \leq x) = 0.10.
	\]
	
	\[
	P(Z<-1.28)=0.10.
	\]
	
	\StandardNormalLeftTailGraph{-1.28}{z=-1.28}{P(Z<-1.28)=0.10}
	
	Convert to standard normal:
	\[
	P\left(Z \leq \frac{x - 38}{7}\right) = 0.10.
	\]
	
	Inverse normal value: $z_{0.10} \approx -1.28$. Thus
	\[
	\frac{x - 38}{7} = -1.28 \Rightarrow x = 38 - 1.28(7) = 29.04.
	\]
	
	\[
	x \approx 29.0\text{ min}
	\]
	
	\textbf{Interpretation.} About 10\% of commutes are 29.0 minutes or shorter.
	
	\section{Central 95\% Device Battery Life}
	
	\textbf{Problem.} Battery life (hours) for a device model follows
	\[
	X \sim N(20, 2.5^2).
	\]
	
	\textbf{Questions/Tasks.} Identify the target probabilities before using inverse normal, then find the interval $[a,b]$ containing the middle 95\% of battery lives.
	
	\subsection*{C7}
	
	Middle 95\% leaves 2.5\% in each tail:
	\[
	P(X \leq a) = 0.025, \qquad P(X \leq b) = 0.975.
	\]
	
	\[
	P(-1.96<Z<1.96)=0.95.
	\]
	
	\StandardNormalCentralBandGraph{-1.96}{1.96}{z=-1.96}{z=1.96}{P(-1.96<Z<1.96)=0.95}
	
	Standardize each endpoint:
	\[
	\frac{a - 20}{2.5} = z_{0.025}, \qquad \frac{b - 20}{2.5} = z_{0.975}.
	\]
	
	From inverse normal values, $z_{0.025} \approx -1.96$, $z_{0.975} \approx 1.96$. So
	\[
	a = 20 - 1.96(2.5) = 15.10, \qquad b = 20 + 1.96(2.5) = 24.90.
	\]
	
	\[
	[a,b] \approx [15.1, 24.9]\text{ h}
	\]
	
	\textbf{Interpretation.} About 95\% of batteries last between 15.1 h and 24.9 h.
	
	\section{Bottom 20\% Streaming Speeds}
	
	\textbf{Problem.} Home streaming speed (Mbps) is modeled by
	\[
	X \sim N(85, 12^2).
	\]
	
	\textbf{Questions/Tasks.} Identify the target probability before using inverse normal, then find the speed $x$ that separates the bottom 20\% of connections.
	
	\subsection*{C7}
	
	Bottom 20\% means
	\[
	P(X \leq x) = 0.20.
	\]
	
	\[
	P(Z<-0.84)=0.20.
	\]
	
	\StandardNormalLeftTailGraph{-0.84}{z=-0.84}{P(Z<-0.84)=0.20}
	
	Convert to $Z$:
	\[
	P\left(Z \leq \frac{x - 85}{12}\right) = 0.20.
	\]
	
	Inverse normal value: $z_{0.20} \approx -0.84$. Then
	\[
	\frac{x - 85}{12} = -0.84 \Rightarrow x = 85 - 0.84(12) = 74.92.
	\]
	
	\[
	x \approx 74.9\text{ Mbps}
	\]
	
	\textbf{Interpretation.} Roughly 20\% of homes have streaming speed below about 74.9 Mbps.
	
	\section{Top 2\% Plant Heights}
	
	\textbf{Problem.} The height of plants in a greenhouse (cm) follows
	\[
	X \sim N(48, 6^2).
	\]
	
	\textbf{Questions/Tasks.} Identify the target probability before using inverse normal, then find the height $x$ so that only the tallest 2\% are above it.
	
	\subsection*{C7}
	
	Top 2\% means
	\[
	P(X > x) = 0.02 \Rightarrow P(X \leq x) = 0.98.
	\]
	
	\[
	P(Z>2.05)=0.02.
	\]
	
	\StandardNormalRightTailGraph{2.05}{z=2.05}{P(Z>2.05)=0.02}
	
	Standardize:
	\[
	P\left(Z \leq \frac{x - 48}{6}\right) = 0.98.
	\]
	
	From the inverse normal table, $z_{0.98} \approx 2.05$. So
	\[
	\frac{x - 48}{6} = 2.05 \Rightarrow x = 48 + 2.05(6) = 60.30.
	\]
	
	\[
	x \approx 60.3\text{ cm}
	\]
	
	\textbf{Interpretation.} Plants taller than about 60.3 cm are in the tallest 2\%.
	
	\section{Bottom 1\% Package Delivery Times}
	
	\textbf{Problem.} Delivery times for a service (hours) are modeled by
	\[
	X \sim N(36, 5^2).
	\]
	
	\textbf{Questions/Tasks.} Identify the target probability before using inverse normal, then find the delivery time $x$ below which the fastest 1\% of deliveries occur.
	
	\subsection*{C7}
	
	Bottom 1\% means
	\[
	P(X \leq x) = 0.01.
	\]
	
	\[
	P(Z<-2.33)=0.01.
	\]
	
	\StandardNormalLeftTailGraph{-2.33}{z=-2.33}{P(Z<-2.33)=0.01}
	
	Convert to a standard normal statement:
	\[
	P\left(Z \leq \frac{x - 36}{5}\right) = 0.01.
	\]
	
	Inverse normal value: $z_{0.01} \approx -2.33$. Then
	\[
	\frac{x - 36}{5} = -2.33 \Rightarrow x = 36 - 2.33(5) = 24.35.
	\]
	
	\[
	x \approx 24.4\text{ h}
	\]
	
	\textbf{Interpretation.} About 1\% of packages are delivered in 24.4 hours or less.
	
	\section{Middle 80\% Fitness Run Times}
	
	\textbf{Problem.} A 2 km run time (minutes) for students is modeled by
	\[
	X \sim N(11.8, 1.4^2).
	\]
	
	\textbf{Questions/Tasks.} Identify the target probabilities before using inverse normal, then find the interval $[a,b]$ containing the middle 80\% of run times.
	
	\subsection*{C7}
	
	Middle 80\% leaves 10\% in each tail:
	\[
	P(X \leq a) = 0.10, \qquad P(X \leq b) = 0.90.
	\]
	
	\[
	P(-1.28<Z<1.28)=0.80.
	\]
	
	\StandardNormalCentralBandGraph{-1.28}{1.28}{z=-1.28}{z=1.28}{P(-1.28<Z<1.28)=0.80}
	
	Standardizing gives
	\[
	\frac{a - 11.8}{1.4} = z_{0.10}, \qquad \frac{b - 11.8}{1.4} = z_{0.90}.
	\]
	
	From inverse normal values: $z_{0.10} \approx -1.28$, $z_{0.90} \approx 1.28$. Hence
	\[
	a = 11.8 - 1.28(1.4) = 10.008, \qquad b = 11.8 + 1.28(1.4) = 13.592.
	\]
	
	\[
	[a,b] \approx [10.0, 13.6]\text{ min}
	\]
	
	\textbf{Interpretation.} About 80\% of students finish the 2 km run between 10.0 and 13.6 minutes.
	
	\section{Top 10\% Product Lifetime}
	
	\textbf{Problem.} A light bulb lifetime (hours) is modeled by
	\[
	X \sim N(1200, 150^2).
	\]
	
	\textbf{Questions/Tasks.} Identify the target probability before using inverse normal, then find the lifetime $x$ that marks the top 10\% longest-lasting bulbs.
	
	\subsection*{C7}
	
	Top 10\% means
	\[
	P(X > x) = 0.10 \Rightarrow P(X \leq x) = 0.90.
	\]
	
	\[
	P(Z>1.28)=0.10.
	\]
	
	\StandardNormalRightTailGraph{1.28}{z=1.28}{P(Z>1.28)=0.10}
	
	Convert to standard normal:
	\[
	P\left(Z \leq \frac{x - 1200}{150}\right) = 0.90.
	\]
	
	Inverse value: $z_{0.90} \approx 1.28$. So
	\[
	\frac{x - 1200}{150} = 1.28 \Rightarrow x = 1200 + 1.28(150) = 1392.
	\]
	
	\[
	x \approx 1392\text{ h}
	\]
	
	\textbf{Interpretation.} Only about 10\% of bulbs last longer than roughly 1392 hours.
	
	\section{Middle 90\% Manufacturing Tolerance Width}
	
	\textbf{Problem.} A metal rod diameter (mm) is modeled by
	\[
	X \sim N(10.00, 0.08^2).
	\]
	
	\textbf{Questions/Tasks.} Identify the target probabilities before using inverse normal, then find the central 90\% interval $[a,b]$ for acceptable diameters.
	
	\subsection*{C7}
	
	Middle 90\% leaves 5\% in each tail:
	\[
	P(X \leq a) = 0.05, \qquad P(X \leq b) = 0.95.
	\]
	
	\[
	P(-1.645<Z<1.645)=0.90.
	\]
	
	\StandardNormalCentralBandGraph{-1.645}{1.645}{z=-1.645}{z=1.645}{P(-1.645<Z<1.645)=0.90}
	
	Standardize:
	\[
	\frac{a - 10.00}{0.08} = z_{0.05}, \qquad \frac{b - 10.00}{0.08} = z_{0.95}.
	\]
	
	Inverse normal values: $z_{0.05} \approx -1.645$, $z_{0.95} \approx 1.645$. Then
	\[
	a = 10.00 - 1.645(0.08) = 9.8684, \qquad b = 10.00 + 1.645(0.08) = 10.1316.
	\]
	
	\[
	[a,b] \approx [9.868, 10.132]\text{ mm}
	\]
	
	\textbf{Interpretation.} About 90\% of rod diameters are expected between 9.868 mm and 10.132 mm.
	
	\section{Middle 96\% Chemistry Scores}
	
	\textbf{Problem.} Chemistry quiz scores are modeled by
	\[
	X \sim N(64, 10^2).
	\]
	
	\textbf{Questions/Tasks.} Identify the target probabilities before using inverse normal, then find the score interval $[a,b]$ containing the middle 96\% of students.
	
	\subsection*{C7}
	
	Middle 96\% leaves 2\% in each tail:
	\[
	P(X \leq a) = 0.02, \qquad P(X \leq b) = 0.98.
	\]
	
	\[
	P(-2.05<Z<2.05)=0.96.
	\]
	
	\StandardNormalCentralBandGraph{-2.05}{2.05}{z=-2.05}{z=2.05}{P(-2.05<Z<2.05)=0.96}
	
	Convert to standard normal equations:
	\[
	\frac{a - 64}{10} = z_{0.02}, \qquad \frac{b - 64}{10} = z_{0.98}.
	\]
	
	From inverse normal values, $z_{0.02} \approx -2.05$, $z_{0.98} \approx 2.05$. So
	\[
	a = 64 - 2.05(10) = 43.5, \qquad b = 64 + 2.05(10) = 84.5.
	\]
	
	\[
	[a,b] \approx [43.5, 84.5]
	\]
	
	\textbf{Interpretation.} About 96\% of students score between roughly 43.5 and 84.5 points.
	
\end{document}